% Generated draft: requires mathematical and source review.
\documentclass[11pt]{article}
\usepackage[margin=25mm]{geometry}
\usepackage{fontspec}
\setmainfont{Latin Modern Roman}
\usepackage{amsmath,amssymb,mathtools,unicode-math}
\setmathfont{Latin Modern Math}
\usepackage{enumitem,needspace,xurl,hyperref}
\setlength{\parindent}{0pt}
\setlength{\parskip}{5pt}
\setlength{\emergencystretch}{3em}
\newcommand{\N}{\mathbb{N}}
\newcommand{\Z}{\mathbb{Z}}
\newcommand{\Q}{\mathbb{Q}}
\newcommand{\R}{\mathbb{R}}
\newcommand{\C}{\mathbb{C}}
\newcommand{\F}{\mathbb{F}}
\newcommand{\A}{\mathbb{A}}
\newcommand{\PP}{\mathbb{P}}
\newcommand{\E}{\mathbb{E}}
\newcommand{\eps}{\varepsilon}
\newcommand{\dd}{\,\mathrm{d}}
\begin{document}
\section*{UnsolvedMath ID: 30002757}
\textbf{Source problem code:} catalogue entry
\subsection{Definitions and mathematical statement}
A fundamental system for a specified scalar second-order linear differential equation is a pair of linearly independent solutions, so that every solution is their linear combination. A normalization prescribes additional conditions selecting such a pair; the source does not specify the conditions that would make a normalization satisfactory.

For terminology only, consider the conventional real scalar Sturm--Liouville setting
\[
 -(p u')'+q u=z w u
\]
on an interval \(I=(a,b)\), where \(p>0\) and \(w>0\) almost everywhere, \(q\) is real-valued, and \(p^{-1},q,w\in L^1_{\mathrm{loc}}(I)\). For a fixed \(z\in\C\setminus\R\), an endpoint is of limit-circle type if every solution is square-integrable near that endpoint with respect to \(w(x)\,\dd x\); it is of limit-point type if the space of such square-integrable solutions has dimension at most one. These are conditional definitions explaining the terminology, not additional hypotheses supplied by the record: the record does not specify its differential expression, coefficient class, interval, or weight.

\textbf{Source problem.} Formulate satisfactory normalizations of a fundamental system of solutions near endpoints of limit-point or limit-circle type which are not regular singular points. The requested formulation should include examples such as the Lennard--Jones potential
\[
 q(x)=4\epsilon\left(\left(\frac{\sigma}{x}\right)^{12}-\left(\frac{\sigma}{x}\right)^6\right),
 \qquad \epsilon>0,\quad \sigma>0.
\]
As a real-variable formula this is defined for \(x\ne0\); the source does not specify the interval, the side of approach to an endpoint, or whether the intended differential equation is specifically \(-u''+q u=z u\).

The request remains a formulation problem rather than a fully specified existence assertion. In particular, the source gives no precise requirements for uniqueness, endpoint asymptotics, spectral-parameter dependence, or admissible changes of basis, and it does not specify the coefficient regularity or convention underlying the exclusion of regular singular points. These unresolved choices must be supplied before a precise normalization theorem can be stated.
\subsection{Short English statement}
Seek satisfactory ways to normalize a basis of solutions near singular Sturm--Liouville endpoints of limit-point or limit-circle type, including the stated Lennard--Jones example with \(\epsilon,\sigma>0\). The equation, domain, and meaning of satisfactory normalization are insufficiently specified for a definitive theorem statement. The supplied open status requires independent review.
\subsection{Sources}
\begin{itemize}
\item Supplied local UnsolvedMath registry record, permanent ID \(30002757\), source problem code \texttt{catalogue entry}, title: \emph{Canonical Endpoint Normalizations for Singular Sturm--Liouville Problems}; collection: Oberwolfach Reports Open Problems. No individual proposer is identified in the supplied record; the proposed year is recorded as 2015.
\item Supplied report citation: \emph{Spectral Theory and Weyl Functions} (2015), Oberwolfach Reports, pp.~5--86, DOI 10.4171/owr/2015/1. \url{https://doi.org/10.4171/owr/2015/1}. The supplied background describes Oberwolfach Reports as documenting mathematical workshops and their open-problem discussions; no problem-specific page is given.
\item Supplied literature source labeled ``original OWR report'': \url{https://doi.org/10.4171/owr/2015/1}.
\item Supplied literature source labeled ``Kostenko et al., singular Sturm--Liouville boundary triples'': \url{https://arxiv.org/abs/2211.08744}. The bibliographic attribution and relevance are supplied and have not been independently verified.
\item Literature assessment checked on \textbf{2026-08-21}: this is supplied dated curation, not an independently verified assessment. It labels the problem open and reports that the cited connection theory treats regular-singular and other special endpoints, while the literature does not supply the requested satisfactory normalization for general ``LP/N'' endpoints such as the Lennard--Jones potential. The supplied background explicitly calls this dated literature triage derived from the OWR curation and requests independent verification of time-sensitive status and cited claims.
\item The registry labels the statement \texttt{corrected\_verified}, explaining that mathematical content was preserved while notation, typography, or omitted definitions were repaired for publication. This is a supplied editorial qualification, not independent mathematical verification of this draft.
\end{itemize}
\end{document}
